\medterm{Xanthelasma} A harmless yellowish patch or bump caused by fat deposits in the skin of the eyelids, usually near the inner corners of the eyes. It can occur with high blood fats but may also occur when blood-fat levels are normal.

\synonyms
Xanthelasma Palpebrarum; Eyelid Xanthoma; Palpebral Xanthoma

\medterm{Xanthine} A natural breakdown product of substances found in cells and foods. It normally changes into uric acid, but can build up when a rare enzyme disorder blocks this pathway and may then contribute to kidney or urinary stones.

\medterm{Xanthine Oxidase} An enzyme that helps convert cell and food breakdown products into uric acid. Medicines used for gout can block this enzyme, lowering uric-acid production and helping prevent urate crystal deposits.

\medterm{Xanthine Oxidase Inhibitors} Medicines such as allopurinol and febuxostat that lower uric acid by blocking an enzyme involved in its production. They are used for selected people with gout or to prevent high uric acid during some cancer treatments. Choice and timing depend on the person’s condition.

\synonyms
XOIs

\medterm{Xanthinuria} A rare inherited disorder in which the body cannot properly change xanthine into uric acid. Xanthine can build up and form kidney or urinary stones; muscle problems may also occur.

\medterm{Xanthochromia} Yellow or amber color in the fluid around the brain and spinal cord. It can occur when blood has broken down in the fluid, including after bleeding around the brain, but other causes are possible and the meaning depends on timing and other test results.

\synonyms
CSF Xanthochromia; Cerebrospinal Fluid Yellowing

\medterm{Xanthogranulomatous} Describing inflammation in which tissue contains many immune cells filled with fat. This pattern can occur in long-lasting infections and other inflammatory conditions.

\medterm{Xanthogranulomatous Cholecystitis} A rare, long-lasting inflammation of the gallbladder that can thicken and damage its wall. It may spread to nearby tissues and can look like gallbladder cancer on scans.

\synonyms
XGC

\medterm{Xanthogranulomatous Oophoritis} A rare long-lasting inflammation of an ovary, often linked to infection or endometriosis. It can form a mass that resembles ovarian cancer on scans and may cause pelvic pain, fever, or abnormal bleeding.

\medterm{Xanthogranulomatous Pyelonephritis} A severe, long-lasting kidney infection that destroys normal kidney tissue and replaces it with inflamed tissue. It is often associated with blockage or stones and may require antibiotics, drainage, or removal of the damaged kidney.

\medterm{Xanthoma} A yellow-orange spot, bump, or patch formed by fat deposits in the skin, tissue beneath the skin, or tendons. It may occur with high blood fats, diabetes, some liver diseases, or an inherited lipid disorder.

\medterm{Xanthoma Disseminatum} A rare disorder in which many yellow, red-brown, or orange bumps develop on the skin, especially in body folds. The mouth, throat, or other organs may also be affected, and some people develop problems controlling water balance.

\synonyms
Montgomery Syndrome; Disseminate Xanthoma

\medterm{Xanthoma Eruptivum} A sudden outbreak of many small yellow-orange bumps, usually on the buttocks, arms, legs, or shoulders. It is often linked to very high blood triglyceride levels.

\synonyms
Eruptive Xanthoma; Eruptive Xanthomatosis

\medterm{Xanthoma Planum} Flat or slightly raised yellow skin patches caused by the buildup of fat inside immune cells. They may occur in skin folds, the palms, or more widely on the body and can be linked to high blood fats or some blood disorders.

\synonyms
Plane Xanthoma; Xanthoma Striatum Palmare

\medterm{Xanthomas} Yellowish papules, plaques, or nodules caused by lipid accumulation in skin immune cells. They may reflect high cholesterol or triglycerides, diabetes, cholestatic liver disease, or an inherited lipid disorder and should prompt appropriate evaluation.
\medterm{Xanthoma Striatum Palmare} Yellow or orange lines or patches in the creases of the palms caused by fat deposits in the skin. They are strongly linked to some inherited blood-fat disorders, especially type III dysbetalipoproteinemia, but blood tests are needed.

\synonyms
Palmar Crease Xanthoma; Planar Palmar Xanthoma

\medterm{Xanthoma Tendinosum} Firm, usually painless fat deposits in tendons or ligaments, most often in the Achilles tendon or tendons of the hands. They are strongly linked to inherited high cholesterol.

\synonyms
Tendon Xanthoma; Tendinous Xanthoma

\medterm{Xanthomatosis} A disorder in which fatty deposits called xanthomas develop in the skin, tendons, or organs. It may be linked to inherited or acquired problems with blood fats.

\medterm{Xanthoma Tuberosum} Slowly enlarging, painless, firm, yellowish-red nodular lipid deposits occurring over extensor pressure areas, particularly the elbows, knees, heels, and buttocks. They are strongly associated with familial hypercholesterolemia and type III hyperlipoproteinemia.

\synonyms
Tuberous Xanthoma; Tuberous Xanthomatosis

\medterm{Xanthopsia} A visual disturbance in which objects appear yellow or yellowish. It can occur with some eye disorders, medicine toxicity such as digoxin toxicity, or disorders of the nerves or retina.

\medterm{Xanthosis} Yellow discoloration of skin or other tissues caused by carotene or another yellow substance. Unlike jaundice, it does not usually make the whites of the eyes yellow; the cause is assessed from the person’s diet and health.

\medterm{X Chromosome} One of the two types of human sex chromosomes. It carries genes involved in sex development as well as many other body functions. Most people have one or two X chromosomes.

\medterm{Xenobiotic} A chemical that does not naturally belong in the body, such as a medicine, pollutant, pesticide, or industrial compound. The body changes and removes it, but some resulting chemicals can be harmful.

\medterm{Xenobiotic Metabolism} The body’s processing of foreign chemicals, including medicines, pollutants, and toxins. Enzymes, especially in the liver, change them so they can be removed; the changed substances may remain active or become harmful.

\medterm{Xenobiotic Resistance} The ability of cells or organisms to withstand the effects of foreign chemicals, medicines, or toxins. It may result from breaking down the substance faster, pumping it out of cells, or changing the substance’s target.

\synonyms
Multidrug Resistance; Xenobiotic Tolerance

\medterm{Xenograft} Tissue, cells, or a biological material transplanted between different species, such as animal tissue used in a human. Rejection and infection are important risks, although some materials are used temporarily as tissue coverings or substitutes.

\medterm{Xenon Anesthesia} Use of xenon gas to cause unconsciousness and prevent pain during an operation. Xenon acts on brain signaling and is cleared mainly through breathing.

\medterm{Xenotransplantation} Transplanting living cells, tissues, or organs between different species, usually from an animal to a human. It may increase the supply of transplant material but carries major risks of immune rejection, infection, and poor organ function.

\medterm{Xerochilia} Abnormal, persistent dryness, scaling, and superficial vertical cracking or fissuring of the lips, resulting from environmental dehydration, persistent lip-licking (cheilitis simplex), chronic mouth breathing, or systemic therapy with oral retinoids (such as isotretinoin).

\synonyms
Dry Lips

\medterm{Xeroderma} Alternate name; see \textit{Dry skin}.

\medterm{Xeroderma Pigmentosum} A rare usually autosomal-recessive disorder of nucleotide-excision DNA repair, with a variant form affecting translesion synthesis. Minimal ultraviolet exposure can cause severe sunburn or persistent redness, early freckle-like pigmentation, and rapid photoaging; skin and ocular-surface cancers occur at a young age, and some people develop progressive neurologic disease. Strict ultraviolet protection, regular skin and eye examinations, and early treatment of suspicious lesions are essential.

\medterm{Xeroderma Pigmentosum Complementation Groups} Genetic subtypes of xeroderma pigmentosum. Each subtype results from a change in a different gene involved in repairing damage caused by ultraviolet light. All subtypes can cause extreme sun sensitivity, freckling, and a very high risk of skin cancer.

\synonyms
XP Complementation Groups; NER Defect Groups

\medterm{Xerophthalmia} Dryness and damage of the eye caused by severe vitamin A deficiency. It may begin with poor night vision and dry eyes and can progress to corneal damage and permanent sight loss.
\synonyms{Vitamin A Deficiency Eye Disease; Ocular Xerosis; Keratomalacia}

\medterm{Xerophthalmic Ulceration} Severe progressive ulceration of the cornea resulting from severe vitamin A deficiency (classified as WHO stage X3A when involving less than one-third of the corneal area, and stage X3B when involving more). Untreated, it rapidly causes corneal liquefactive necrosis (keratomalacia) and permanent blindness.

\synonyms
Vitamin A Corneal Ulcer; Keratomalacic Ulceration

\medterm{Xerosis} Alternate name; see \textit{Dry skin}.

\medterm{Xerostomia} The sensation of dry mouth caused by reduced or altered saliva production. It may result from
medicines, dehydration, salivary-gland disease, radiation, or systemic illness and can impair speaking, swallowing,
taste, and dental health.

\medterm{Xerotic Eczema} A common, intensely pruritic inflammatory dermatosis occurring predominantly on the pretibial lower extremities of elderly individuals during cold, dry winter months, characterized by profound epidermal dehydration, scaling, and fine superficial erythematous fissures resembling "cracked porcelain" or a "dry riverbed."

\synonyms
Asteatotic Eczema; Eczema Craquele; Winter Itch

\medterm{X-Inactivation} A process in early development that switches off most genes on one of two X chromosomes in each cell. It helps balance gene activity, although some genes remain active and the pattern can differ between cells.

\medterm{Xiphisternal Joint} The joint between the lower part of the breastbone and the xiphoid process. It often becomes partly or fully hardened with age.

\synonyms
Xiphisternal Synchondrosis; Xiphisternal Articulation

\medterm{Xiphoidalgia} Pain or tenderness over the xiphoid process, the small lower part of the breastbone. Pain may spread to the upper abdomen, chest, back, or shoulders and may worsen with touch, bending, lifting, or coughing.

\synonyms
Xiphoidodynia; Hypersensitive Xiphoid Syndrome; Xiphoid Syndrome

\medterm{Xiphoidectomy} Surgery to remove the xiphoid process, the small lower part of the breastbone. It may be done for persistent pain in this area or to provide access during some chest or abdominal operations.

\medterm{Xiphoid Process} The small pointed lower part of the breastbone. Its shape varies between people and it provides attachment for some muscles and ligaments.

\medterm{X-Linked} Describing a gene or condition caused by a change on the X chromosome. The pattern differs between people with one X chromosome and those with two; fathers pass their X chromosome to daughters, not sons.

\medterm{X-Linked Adrenoleukodystrophy} An inherited disorder, usually affecting boys, in which certain fats build up and damage the adrenal glands and nervous system. It can cause behavior or learning changes, movement problems, seizures, or adrenal failure.

\medterm{X-Linked Agammaglobulinemia} An inherited immune disorder, usually affecting boys, in which the body cannot make enough antibodies. Repeated bacterial infections often begin in infancy or early childhood and are treated with antibody replacement and infection care.

\medterm{X-Linked Chronic Granulomatous Disease} An inherited immune disorder, usually affecting boys, caused by a change in the \textit{CYBB} gene. Certain white blood cells cannot kill some bacteria and fungi normally, leading to repeated serious infections and areas of ongoing inflammation.

\synonyms
X-Linked CGD; CYBB Deficiency; Cytochrome b-245 Beta Deficiency

\medterm{X-Linked Dominant Inheritance} A pattern in which one changed gene on the X chromosome can cause a condition. An affected father passes the changed X chromosome to all daughters and no sons.

\medterm{X-Linked Emery-Dreifuss Muscular Dystrophy} An inherited muscle disorder, usually affecting boys, that causes stiff joints, slowly worsening muscle weakness, and problems with the heart’s electrical rhythm.

\synonyms
EDMD1; Emerin Deficiency Muscular Dystrophy

\medterm{X-Linked Hyper-IgM Syndrome} An inherited immune disorder, usually affecting boys, in which the body cannot change one type of protective antibody into other types. It causes repeated lung and sinus infections and may cause infections by organisms that usually do not cause illness.

\synonyms
HIGM1; CD40L Deficiency; Hyper-IgM Syndrome Type 1

\medterm{X-Linked Hypophosphatemia} An inherited disorder that causes the kidneys to lose too much phosphate. Low phosphate can lead to soft or poorly growing bones, bone pain, bowed legs, dental problems, and reduced height.

\medterm{X-Linked Ichthyosis} An inherited, usually X-linked recessive skin disorder caused by a variant or deletion affecting the STS gene and deficiency of steroid sulfatase. It usually affects males and causes widespread dry, dark, plate-like scales that often begin in infancy or early childhood.

\medterm{X-Linked Inheritance} Inheritance of a gene change located on the X chromosome. The effects depend on the gene, the type of change, and whether a person has one or two X chromosomes; fathers do not pass an X chromosome to sons.

\medterm{X-Linked Lymphoproliferative Disease Type 1} An inherited immune disorder, usually affecting boys, in which the body can react severely to Epstein--Barr virus. It may cause severe glandular fever, very high inflammation, low antibody levels, or lymphoma.

\synonyms
XLP1; Duncan Disease; SAP Deficiency

\medterm{X-Linked Lymphoproliferative Syndrome} An inherited immune disorder, usually affecting boys, in which the body reacts abnormally to infections. Severe Epstein–Barr virus illness, persistent fever, liver or blood problems, lymphoma, and intestinal inflammation may occur.

\medterm{X-Linked Myotubular Myopathy} A rare inherited muscle disease, usually affecting boys, that causes severe weakness and low muscle tone from birth. Babies may have drooping eyelids, feeding difficulty, and dangerous weakness of the breathing muscles.

\synonyms
XLMTM; Myotubularin Deficiency

\medterm{X-Linked Recessive Inheritance} A pattern in which a changed gene on the X chromosome usually causes a condition in people with one X chromosome. People with two X chromosomes are often unaffected carriers, although they can sometimes have symptoms. An affected father does not pass the changed X chromosome to sons.

\medterm{X-Linked Retinitis Pigmentosa} An inherited disorder in which the retina slowly loses its ability to sense light. It usually causes night blindness followed by narrowing of the field of vision and worsening sight.

\synonyms
XLRP; X-Linked Pigmentary Retinopathy

\medterm{X-Linked Retinoschisis} An inherited eye disorder, usually affecting boys, in which layers of the retina separate. It can cause blurred central or distance vision, crossed eyes, bleeding, or retinal detachment in severe cases.

\medterm{X-Linked Severe Combined Immunodeficiency} A severe inherited immune disorder, usually affecting boys, in which important immune cells do not develop normally and antibody protection is poor. Babies are highly vulnerable to serious infections.

\medterm{X-Linked Sideroblastic Anemia} An inherited form of anemia, usually affecting boys, in which red blood cells cannot make hemoglobin normally. Iron builds up in developing blood cells, causing small, pale red blood cells and anemia.

\synonyms
XLSA; ALAS2 Deficiency; Hereditary Sideroblastic Anemia

\medterm{X-Linked Spinal and Bulbar Muscular Atrophy} An inherited muscle and nerve disorder that usually begins in adulthood and slowly causes weakness. It may also cause muscle twitching, trouble speaking or swallowing, breast enlargement, and reduced fertility.

\synonyms
Kennedy Disease; SBMA; Bulbo-Spinal Muscular Atrophy

\medterm{X-Ray} An imaging test that uses a small amount of radiation to create a picture of structures inside the body. It is especially useful for bones, lungs, teeth, and some other tissues.

\medterm{X-Ray Attenuation} The reduction in strength of an X-ray beam as it passes through the body. Bone, soft tissue, and air reduce the beam by different amounts, creating the differences seen in the image.

\medterm{X-Ray Crystallography} A laboratory method that uses X-rays to work out the three-dimensional structure of molecules such as proteins. The molecule is made into a crystal, and the pattern produced by the X-rays is analyzed.

\synonyms
X-Ray Crystal Diffraction Analysis; Macromolecular Crystallography

\medterm{X-Ray Shield} Protective material, often lead or another dense substance, used to reduce exposure to scattered X-rays during imaging.

\medterm{X-Ray Therapy} Treatment that uses carefully aimed X-ray radiation to damage abnormal cells or relieve symptoms. It is used mainly for cancer and may be given from a machine outside the body or near the target.

\medterm{XX Male Syndrome} A difference in sex development in which a person with two X chromosomes develops typically male physical features, usually because a sex-development gene has moved onto an X chromosome. Infertility is common, and care is individualized.

\medterm{XXX Syndrome} Alternate name; see \textit{Triple X syndrome}, a sex-chromosome variation in which a female has an additional X chromosome.

\medterm{XXXX Syndrome} XXXX syndrome is a rare sex-chromosome variation in which a female has four X chromosomes; it may cause developmental or learning difficulties, distinctive physical features, and ovarian dysfunction, with severity varying widely.

\medterm{Xylitol} A sugar alcohol used as a sweetener and in some oral-care products. Large human ingestions may cause gastrointestinal symptoms; it is highly toxic to dogs.

\medterm{Xylose} A simple sugar found in some plants. A form called D-xylose has been used in tests of how well the small intestine absorbs nutrients and is also used in laboratory studies.

\medterm{Xylose Absorption Test} A test of how well the small intestine absorbs a simple sugar called D-xylose. Low amounts in the blood or urine may suggest poor absorption from the intestine; normal results with other signs of poor absorption may suggest a problem with the pancreas.

\synonyms
D-Xylose Absorption Test; D-Xylose Tolerance Test

\medterm{XYY Syndrome} A sex-chromosome variation in which a male has an extra Y chromosome. Many affected individuals have few or no obvious symptoms, though some have tall stature, learning difficulties, or developmental differences.

\medterm{X-linked Juvenile Retinoschisis} An inherited retinal disorder, usually affecting boys, in which layers of the retina separate during childhood. It can cause reduced central vision, difficulty seeing in dim light, crossed eyes, or retinal complications.

\medterm{X-linked Adult Spinal Muscular Atrophy} Alternate name for Kennedy disease, or X-linked spinal and bulbar muscular atrophy. It is an inherited androgen-receptor disorder that gradually causes weakness, muscle wasting, twitching, swallowing or speech difficulty, and sometimes breast enlargement.

\medterm{X-linked Intellectual Disability} Intellectual disability caused by a change in a gene on the X chromosome. Features vary widely and may include developmental delay, learning difficulty, behavior differences, seizures, or associated physical findings.

\medterm{XO Syndrome} A form of Turner syndrome caused by complete or partial absence of one X chromosome. It may cause short stature, ovarian insufficiency, infertility, characteristic physical findings, and heart, kidney, hearing, or thyroid problems.

\medterm{XP} Abbreviation for xeroderma pigmentosum, a rare DNA-repair disorder causing extreme sensitivity to ultraviolet light and a markedly increased risk of early skin and eye cancers.

\medterm{XXY Males} People with a typically male pattern of development and an extra X chromosome, usually 47,XXY. This is commonly called Klinefelter syndrome and may cause small testes, low testosterone, infertility, tall stature, or learning and language difficulties.

\medterm{XXY Syndrome} Alternate name for Klinefelter syndrome, a sex-chromosome variation in which a typically male individual has an extra X chromosome. Effects vary and may include low testosterone, infertility, small testes, tall stature, and developmental or language differences.
